\documentclass[]{book}
\usepackage{fontspec}
\usepackage{xunicode}
\usepackage{polyglossia}
\setmainlanguage{french}
\usepackage{lipsum}
%Titre, auteur et date se déclarent dans le préambule...
\title{Structurer son document}
\author{Prénom Nom}
\date{4 novembre 2026}

\begin{document}
%Comment structurer son document

%faire apparaître le titre 
\maketitle%... et s'impriment dans le corps du document avec \maketitle

%Mettre un titre de chapitre ici
\chapter{Un chapitre}%\chapter n'existe que dans la classe book (et report)

%Mettre un titre de section
\section{Une section}

\lipsum[1-2]
%Mettre un titre de sous-section
\subsection{Une sous-section}

\lipsum[3-4]
%Section
\section{Une autre section}

%sous-section
\subsection{Une autre sous-section}

%sous-sous-section
\subsubsection{Une sous-sous-section}%non numérotée et absente de la table des matières dans la classe book

\lipsum[5-6]

%Imprimer la table des matières
\tableofcontents
%Il faut compiler deux fois : à la première compilation, les titres sont écrits dans le fichier .toc ;
%à la seconde, \tableofcontents lit ce fichier et imprime la table.

\end{document}
